\documentclass{article}
\usepackage{fullpage}
\usepackage{hyperref}
\providecommand{\e}[1]{\ensuremath{\times 10^{#1}}}
\begin{document}

\title{Empirical IO I: Problem Set 0} 
\author{Chris Conlon}
\date{}
\maketitle
This problem set is designed to make sure that your $\mathtt{numpy}$ skills are up to speed. Python/$\mathtt{numpy}$ is the reference implementation below; any language is acceptable, but then it is on you to map each task to the corresponding routines. For each question, the expectation is that you complete the task, provide the appropriate code, and fully read the documentation, and work the $\mathtt{numpy}$ provided examples on your own.  The tasks should be fairly easy.  The goal is to accomplish the tasks in the most straightforward manner with the least amount of code.  

If you are new to $\mathtt{numpy}$ I suggest the following tutorials:
\begin{itemize}
\item \url{https://numpy.org/doc/stable/user/quickstart.html}
\item  \url{https://numpy.org/doc/stable/user/numpy-for-matlab-users.html}
\end{itemize}

\section*{Part 0: Logit Inclusive Value}
The logit inclusive value or $IV = \log \sum_{i=0}^N \exp[x_i]$.
\begin{enumerate}
\item Show that this function is everywhere convex.
\item A common problem in practice is that if one of the $x_i > 600$ that we have an ``overflow'' error on a computer. In this case $\exp[600] \approx 10^{260}$ which is too large to store with any real precision, especially if another $x$ has a different scale (say $x_2=10$). A common ``trick'' is to subtract off $m_i = \max_i x_i$ from all $x_i$.  Show how to implement the trick and get the correct value of $IV$. If you get stuck take a look at Wikipedia.
\item Compare your function to $\mathtt{scipy.special.logsumexp}$. Does it appear to suffer from underflow/overflow? Does it use the $\max$ trick?
\end{enumerate}


\section*{Part 1: Markov Chains}
Consider the following Markov TPM:\\
Let $P = \{p_{i,j} \}$ be an $n \times n$ transition matrix of a Markov process where $\{p_{i,j} \}$ is interpreted as the probability that the system, when it is in state $i$, will move to state $j$. If we denote by $\mu_t = [ \mu_{t,1} , \mu_{t,2}, \ldots, \mu_{t,n}]$ the probability mass function of the system over the $n$ states then $\mu_{t,j}$ evolves according to
\begin{eqnarray*}
\mu_{t+1,j} = \sum_{i=1}^n p_{i,j} \mu_{t,i}
\end{eqnarray*}
Then we can write the state to state transition matrix as :
\begin{eqnarray*}
\mu_{t+1}& &= \mu_t P \\
P &=&
\left[ {\begin{array}{ccc}
    0.2&    0.4&    0.4\\
    0.1&    0.3&    0.6\\
    0.5&    0.1&    0.4\\
 \end{array} } \right]
\end{eqnarray*}
We're interested in the ergodic distribution $\mu P = \mu$. This is similar to the transition matrix infinitely many periods into the future $P^{\infty}$. Write a function that computes the ergodic distribution of the matrix $P$ by examining the properly rescaled eigenvectors and compare your result to $P^{100}$. Here I recommend $\mathtt{numpy.linalg.matrix\_power}$ and $\mathtt{numpy.linalg.eig}$. (A common mistake is element-wise exponentiation of the matrix).


\section*{Part 2: Numerical Integration}
\textit{Note: The \texttt{scipy} quadrature routines may return a set of nodes/weights that correspond to integrating $e^{-x^2}$ over $(\infty,\infty)$, while the nodes/weights available at \url{https://www.sparse-grids.de} may not. It is always important to make sure you understand what your nodes/weights correspond to. One way to do this is to integrate some simple functions $f(x)=1,f(x)=x,f(x)=x^2$ where you know the analytic result and see what the quadrature routine gives as the answer.}\\

In this part we will look at calculating the logit choice probability $\sigma(X,\theta)$ by numerical integration:\\
 $\sigma(X,\theta) =\int_{-\infty}^{\infty} \frac{\exp(\beta_i X)}{1+ \exp(\beta_i X)} f(\beta_i | \theta) \partial \beta_i$. \\
 Assume $f \sim N(0.5,2)$ and that $X = 0.5$.
 
\begin{enumerate}
\item Create the function in an Python called $\mathtt{binomiallogit}$. (It should take $\beta$ the item you integrate over as its argument, it should take the PDF $\mathtt{scipy.stats.norm.pdf}$ as an optional argument).

\item Integrate the function using Python's $\mathtt{scipy.integrate.quad}$ command and setting the tolerance to $1\e{-14}$.  Treat this a the ``true'' value.

\item Integrate the function by taking 20 and 400 Monte Carlo draws from $f$ and computing the sample mean.
\item Integrate the function using Gauss-Hermite quadrature for $k=4, 12$ (Try some odd ones too). Obtain the quadrature nodes and weights from $\mathtt{numpy.polynomial.hermite.hermgauss}$. Gauss-Hermite quadrature assumes a weighting function of $\exp[-x^2]$, you will need a change of variables to integrate over a normal density (see the Numerical Integration notes in \texttt{Extra Notes/}). You also need to pay attention to the constant of integration.
\item Compare results to the Monte Carlo results. \textit{Make sure your quadrature weights sum to 1!}

\item Repeat the exercise in two dimensions where $\mu = (0.5,1), \sigma = (2,1)$, and $X=(0.5,1)$.
\item Now the \textbf{curse of dimensionality}: repeat the exercise in \textbf{four} dimensions with $\mu = (0.5, 1, 0, -0.5)$, $\sigma = (1, 0.5, 0.75, 0.5)$, and $X = (0.5, 1, -1, 0.5)$, comparing:
\begin{enumerate}
\item Smart check first: show that because $\beta' X = \sum_k \beta_k X_k$ is a linear combination of independent normals, it is itself univariate $N\left(\mu'X, \sum_k (\sigma_k X_k)^2\right)$ --- so this ``4-D'' integral is secretly 1-D. Compute the \textbf{true} value with $\mathtt{scipy.integrate.quad}$ on the reduced 1-D integral (tolerance $1\e{-14}$). Lesson: understand the structure of your integrand before reaching for brute force.
\item $\mathtt{scipy.integrate.nquad}$ on the full 4-D integral at a \emph{moderate} tolerance (\texttt{epsabs=1e-6}): \textbf{time it} and report the number of function evaluations (\texttt{full\_output=1}). Adaptive quadrature does not scale; that is part of the lesson. (Do not be tempted to tighten the tolerance --- budget your afternoon.)
\item Tensor-product Gauss-Hermite with $k=5$ and $k=9$ nodes per dimension ($5^4 = 625$ and $9^4=6561$ points).
\item A \textbf{sparse grid} at two or three accuracy levels. Either download the precomputed Gaussian-kernel (GQN) nodes and weights from \url{https://www.sparse-grids.de} (Heiss and Winschel 2008), or construct the grid yourself from 1-D $\mathtt{hermgauss}$ rules via the Smolyak combination technique. \textbf{Required check either way}: before using any rule, verify it on monomials --- $f=1$, $f = x_k$, $f = x_k^2$, $f=x_1^2 x_2^2$ --- where you know the answer under the Gaussian weight. (This is the point of the note above: nodes and weights are only as good as your understanding of the kernel they assume.)
\item Monte Carlo with the same number of draws as each rule above uses nodes.
\item Tensor-product Gauss-Hermite probably looks surprisingly good in your table. Two sentences: why does it do so well on \emph{this} integrand at $d=4$ (hint: your answer to (a)), and why will it lose anyway at $d=10$? Report the node counts for $k=5$ per dimension at $d=10$ versus a level-5 sparse grid.
\end{enumerate}
\item Put everything into three tables (1-D, 2-D, 4-D), showing the error from the ``true'' value and the number of points used in each evaluation. In the 4-D table, which method would you use inside an estimation routine that must evaluate this integral thousands of times?
\item Now Construct a new function $\mathtt{binomiallogitmixture}$ that takes a vector for $X$ and returns a vector of binomial probabilities (appropriately integrated over $f(\beta_i | \theta)$ for the 1-D mixture). It should be obvious that Gauss-Hermite is the most efficient way to do this. \textit{Do NOT use loops}
\end{enumerate}

\section*{Part 3: Fixed Points and the Berry Inversion}
Consider a three-product logit market plus an outside good, with model shares
\begin{eqnarray*}
\sigma_j(\delta) = \frac{\exp[\delta_j]}{1+\sum_{k=1}^{3}\exp[\delta_k]}, \quad j = 1,2,3
\end{eqnarray*}
and observed shares $s = (0.30, 0.20, 0.15)$ (so $s_0 = 0.35$).
\begin{enumerate}
\item Solve for $\delta$ analytically: $\delta_j = \log s_j - \log s_0$ (the Berry inversion from lecture). Verify $\sigma(\delta) = s$.
\item Now pretend you didn't know the closed form. Implement the fixed-point iteration
$\delta^{(t+1)} = \delta^{(t)} + \log s - \log \sigma(\delta^{(t)})$
starting from $\delta^{(0)} = 0$. How many iterations to reach $\|\delta^{(t+1)}-\delta^{(t)}\|_{\infty} < 10^{-13}$? (Use your Part 0 $\mathtt{logsumexp}$ machinery to compute $\log \sigma$ safely.)
\item Solve the same system with a Newton-type solver ($\mathtt{scipy.optimize.root}$) and compare iteration counts and wall-clock time.
\item One sentence: where will you see this fixed point again in this course? (Hint: twice --- once in demand estimation, once in dynamics.)
\end{enumerate}

\section*{Part 4: Numerical Derivatives}
Using your 1-D $\mathtt{binomiallogit}$ mixture $\sigma(X,\theta)$ from Part 2 (Gauss-Hermite version, $k=12$), we want $\frac{\partial \sigma}{\partial X}$ at $X=0.5$.
\begin{enumerate}
\item Compute the derivative by forward differences $\frac{f(X+h)-f(X)}{h}$ and central differences $\frac{f(X+h)-f(X-h)}{2h}$ for $h \in \{10^{-1}, 10^{-2}, \ldots, 10^{-12}\}$.
\item Compute the ``exact'' answer by differentiating under the integral sign (the integrand's derivative has a closed form --- write it down) and applying the same quadrature rule.
\item Plot the absolute error of both difference schemes against $h$ on a log-log scale. Explain the two regimes you see: where does truncation error dominate, and where does floating-point roundoff dominate? What are the (approximate) optimal $h$'s, and how do they relate to machine epsilon?
\item Why does this matter? You are about to spend a semester handing gradients to $\mathtt{scipy.optimize}$. One sentence on when you should prefer analytic (or automatic) derivatives to finite differences.
\end{enumerate}

\end{document}

